%------------------------------------------------------------------------------%
\documentclass[fleqn,utf8,aspectratio=169,12pt]{beamer}

\usepackage{calculo_varias_variaveis-1}

\title{Diferenciais}

%------------------------------------------------------------------------------%
\begin{document}

\addtitle

%------------------------------------------------------------------------------%
\section{Variação em uma Direção}

%------------------------------------------------------------------------------%
\begin{frame}{Variação em uma Dimensão}
  Em uma dimensão podemos estimar a variação do valor da função
  em um ponto próximo à $a$ usando a derivada em $a$

  \pause
  \[
    df = f'(a) dx
  \]

  \[
    \pause
    dx = x - a
  \]

  \[
    \pause
    df \approx f(x) - f(a)
  \]
\end{frame}

%------------------------------------------------------------------------------%
\begin{frame}{Variação em Várias Dimensões}
  Em várias dimensões usamos a derivada direcional
  \pause
  \[
    df
    = \left(D_u f(a, b){}\right) ds
    \pause
    = \left(\nabla f(a, b){} \cdot u\right) ds
  \]
  \pause
  \[
    ds = \sqrt{(x - a)^2 + (y-b)^2}
  \]
  \pause
  \[
    df \approx f(x, y) - f(a, b)
  \]
\end{frame}

%------------------------------------------------------------------------------%
%------------------------------------------------------------------------------%
\begin{question}
  Calcule quanto
  \[
    f(x, y, z) = y\sin(x) + 2yz
  \]
  irá variar se nos movermos 0,1 unidades a partir de \((0, 1, 0)\) \\
  na direção do ponto \((2, 2, -2)\)
\end{question}

%------------------------------------------------------------------------------%
\begin{resolution}{Direção do movimento}
  \[
    \pause
    v
    \pause
    \;=\; \vetortri{2}{2}{-2}
        - \vetortri{0}{1}{ 0}
    \pause
    \;=\; \vetortri{2}{1}{-2}
    \pause
    \;=\; 2\veci + \vecj -2\veck
  \]
  \pause
  Normalizando o vetor
  \[
    \pause
    u
    \pause
    = \frac{v}{\abs{v}}
    \pause
    = \frac{2\veci + \vecj -2\veck}{\sqrt{2^2 + 1^2 + (-2)^2}}
    \pause
    = \frac{2\veci + \vecj -2\veck}{\sqrt{9}}
    \pause
    = \frac{2}{3} \veci + \frac{1}{3} \vecj - \frac{2}{3} \veck
  \]
\end{resolution}

%------------------------------------------------------------------------------%
\begin{resolution}{Derivadas parciais}
  \begin{align*}
    \onslide<+->{\D{f}{x}}
    \onslide<+->{& = \D{}{x}\left(y\sin(x) + 2yz\right)}
    \onslide<+->{ = y\cos(x) \\[5mm]}
    \onslide<+->{\D{f}{y} & = \D{}{y}\left(y\sin(x) + 2yz\right)}
    \onslide<+->{ = \sin(x) + 2z\\[5mm]}
    \onslide<+->{\D{f}{z} & = \D{}{z}\left(y\sin(x) + 2yz\right)}
    \onslide<+->{ = 2y}
  \end{align*}
\end{resolution}

%------------------------------------------------------------------------------%
\begin{resolution}{Gradiente de $f$}

  \[
    \onslide<+->{\nabla f}
    \onslide<+->{= (y\cos(x))\veci + (\sin(x) + 2z)\vecj + (2y)\veck}
    \onslide<+->{= \vetortri{y\cos(x)}{\sin(x) + 2z}{2y}}
  \]

  \[
    \onslide<+->{\nabla f(0, 1, 0)}
    \onslide<+->{= (1\cos(0))\veci + (\sin(0) + 2\times 0)\vecj + (2\times 1)\veck}
    \onslide<+->{= \veci + 2\veck}
    \onslide<+->{= \vetortri{1}{0}{2}}
  \]
\end{resolution}

%------------------------------------------------------------------------------%
\begin{resolution}{Derivada direcional}
  \[
    D_u f(0, 1, 0)
    \pause
    \;=\; \nabla f(0, 1, 0) \cdot u
  \]

  \[
    \hphantom{D_u f(0, 1, 0)}
    \pause
    \;=\; \left(\veci + 2\veck\right)
    \cdot
    \left(\frac{2}{3} \veci + \frac{1}{3} \vecj - \frac{2}{3} \veck\right)
  \]

  \[
    \hphantom{D_u f(0, 1, 0)}
    \pause
    \;=\; 1\times\frac{2}{3} + 0\times\frac{1}{3} + 2 \times\left(- \frac{2}{3}\right)
  \]

  \[
    \hphantom{D_u f(0, 1, 0)}
    \pause
    \;=\; \frac{2}{3} - \frac{4}{3}
    \pause
    \;=\; -\frac{2}{3}
  \]
\end{resolution}

%------------------------------------------------------------------------------%
\begin{resolution}{Variação}
  \[
    \pause
    df
    \pause
    \;=\; D_u f(0, 1, 0) ds
  \]

  \[
    \pause
    \hphantom{df}
    \;=\; -\frac{2}{3}\times 0,1
    \pause
    \;=\; -\frac{2}{3} \frac{1}{10}
    \pause
    \;=\; -\frac{2}{30}
    \pause
    \;=\; -\frac{1}{15}
  \]

  \[
    \pause
    \hphantom{df}
    \approx -0,067
  \]
\end{resolution}

%------------------------------------------------------------------------------%

%------------------------------------------------------------------------------%
\section{Diferenciais}

%------------------------------------------------------------------------------%
\begin{frame}{Diferencial em Uma Dimensão}
  Quando $x$ varia de $a$ para
  \tabto{55mm}\(a+\Delta x\)

  \pause
  \vspace{\baselineskip}
  A variação da função $f$ é
  \tabto{55mm}\(\Delta f = f(a + \Delta x) - f(a)\)

  \pause
  \vspace{\baselineskip}
  A diferencial de $f$ é
  \tabto{55mm}\(df = f'(a) \Delta x\)
\end{frame}

%------------------------------------------------------------------------------%
\begin{frame}{Diferencial em Uma Dimensão}
  Quando $x$ varia de $a$ para
  \tabto{55mm}\(a+dx\)

  \vspace{\baselineskip}
  A variação da função $f$ é
  \tabto{55mm}\(\Delta f = f(a + dx) - f(a)\)

  \vspace{\baselineskip}
  A diferencial de $f$ é
  \tabto{55mm}\(df = f'(a) dx\)
\end{frame}

%------------------------------------------------------------------------------%
\begin{frame}{Diferencial em Duas Dimensão}
  Quando $(x, y)$ varia de $(a, b)$ para
  \tabto{65mm}\((a+dx, b+dy)\)

  \pause
  \vspace{\baselineskip}
  A variação da função $f$ é
  \tabto{65mm}\(\Delta f = f(a+dx, b+dy) - f(a, b)\)

  \pause
  \vspace{\baselineskip}
  A diferencial de $f$ é
  \tabto{65mm}\(df = f_x(a, b) dx +  f_y(a, b) dy\)

  \pause
  \vspace{\baselineskip}
  As diferenciais $dx$ e $dy$ são variáveis independentes
  \[
    dx = x - a
    \qquad
    dy=y-b
  \]
\end{frame}

%------------------------------------------------------------------------------%
\begin{frame}{Definição}
  Se movemos de \((a, b)\) para um ponto próximo \((a+dx, b+dy)\)

  \pause
  \vspace{\baselineskip}
  A \emph{Diferencial Total de $f$} é
  \[
    df \;=\; f_x(a, b)\, dx \;+\;  f_y(a, b)\, dy
  \]
\end{frame}

%------------------------------------------------------------------------------%
%------------------------------------------------------------------------------%
\begin{question}
  Uma lata de formato cilíndrico foi projetada para ter um raio de 1 pol.
  e uma altura de 5 pol., mas o raio tem um erro de \(dr=0,03\) e a altura
  tem um erro de \(dh=-0,1\).

  \vspace{\baselineskip}
  Qual a variação resultante no volume da lata?
\end{question}

\begin{resolution}{Variação exata}
  \begin{align*}
    \onslide<+->{\Delta V}
    \onslide<+->{& = V(r_0+dr, h_0+dh)- V(r_0, h_0) \\}
    \onslide<+->{& = V(1+0,03, 5-0,1) - V(1, 5)\\}
    \onslide<+->{& = V(1,03, 4,9) - V(1, 5) \\}
    \onslide<+->{& = \pi(1,03)^2 \times 4,9 - \pi 1^2 \times 5 \\}
    \onslide<+->{& = 16,3313 - 15,7080 \\}
    \onslide<+->{& = 0,6233}
  \end{align*}
\end{resolution}

\begin{resolution}{Variação usando o diferencial}
  \begin{align*}
    \onslide<+->{\Delta V}
    \onslide<+->{& \approx dV = V_r(r_0, h_0) dr + V_h(r_0, h_0) dh \\}
    \onslide<+->{& = \left(2\pi r h\right)\evalat{(r_0, h_0)}{} dr
      + \left(\pi r^2\right)\evalat{(r_0, h_0)}{} dh \\}
    \onslide<+->{& = \left(2\pi r h\right)\evalat{(1, 5)}{} (0,03)
      + \left(\pi r^2\right)\evalat{(1, 5)}{} (-0,1) \\}
    \onslide<+->{& = \left(2\pi 1 \times 5\right)(0,03)
      - \left(\pi 1^2\right)(0,1) \\}
    \onslide<+->{& = 0,3\pi - 0,1\pi  \\}
    \onslide<+->{& = 0,2\pi \\}
    \onslide<+->{& \approx 0,6283}
    \onslide<+->{& & \Delta V = 0,6233}
  \end{align*}
\end{resolution}

%------------------------------------------------------------------------------%


%------------------------------------------------------------------------------%
\section{Lista Mínima}

%------------------------------------------------------------------------------%
\begin{frame}{Lista Mínima}
  Cálculo Vol. 2 do Thomas 12\textsuperscript{a} ed. -- Seção 14.6
  \begin{enumerate}
    \item Estudar o texto da seção
    \item Resolver os exercícios: 19-24
  \end{enumerate}

  \vfill
  Atenção: A prova é baseada no livro, não nas apresentações
\end{frame}

%------------------------------------------------------------------------------%
\end{document}
%------------------------------------------------------------------------------%
